% Options for packages loaded elsewhere
\PassOptionsToPackage{unicode}{hyperref}
\PassOptionsToPackage{hyphens}{url}
\PassOptionsToPackage{dvipsnames,svgnames,x11names}{xcolor}
%
\documentclass[
  letterpaper,
  DIV=11,
  numbers=noendperiod]{scrartcl}

\usepackage{amsmath,amssymb}
\usepackage{iftex}
\ifPDFTeX
  \usepackage[T1]{fontenc}
  \usepackage[utf8]{inputenc}
  \usepackage{textcomp} % provide euro and other symbols
\else % if luatex or xetex
  \usepackage{unicode-math}
  \defaultfontfeatures{Scale=MatchLowercase}
  \defaultfontfeatures[\rmfamily]{Ligatures=TeX,Scale=1}
\fi
\usepackage{lmodern}
\ifPDFTeX\else  
    % xetex/luatex font selection
\fi
% Use upquote if available, for straight quotes in verbatim environments
\IfFileExists{upquote.sty}{\usepackage{upquote}}{}
\IfFileExists{microtype.sty}{% use microtype if available
  \usepackage[]{microtype}
  \UseMicrotypeSet[protrusion]{basicmath} % disable protrusion for tt fonts
}{}
\makeatletter
\@ifundefined{KOMAClassName}{% if non-KOMA class
  \IfFileExists{parskip.sty}{%
    \usepackage{parskip}
  }{% else
    \setlength{\parindent}{0pt}
    \setlength{\parskip}{6pt plus 2pt minus 1pt}}
}{% if KOMA class
  \KOMAoptions{parskip=half}}
\makeatother
\usepackage{xcolor}
\setlength{\emergencystretch}{3em} % prevent overfull lines
\setcounter{secnumdepth}{-\maxdimen} % remove section numbering
% Make \paragraph and \subparagraph free-standing
\makeatletter
\ifx\paragraph\undefined\else
  \let\oldparagraph\paragraph
  \renewcommand{\paragraph}{
    \@ifstar
      \xxxParagraphStar
      \xxxParagraphNoStar
  }
  \newcommand{\xxxParagraphStar}[1]{\oldparagraph*{#1}\mbox{}}
  \newcommand{\xxxParagraphNoStar}[1]{\oldparagraph{#1}\mbox{}}
\fi
\ifx\subparagraph\undefined\else
  \let\oldsubparagraph\subparagraph
  \renewcommand{\subparagraph}{
    \@ifstar
      \xxxSubParagraphStar
      \xxxSubParagraphNoStar
  }
  \newcommand{\xxxSubParagraphStar}[1]{\oldsubparagraph*{#1}\mbox{}}
  \newcommand{\xxxSubParagraphNoStar}[1]{\oldsubparagraph{#1}\mbox{}}
\fi
\makeatother


\providecommand{\tightlist}{%
  \setlength{\itemsep}{0pt}\setlength{\parskip}{0pt}}\usepackage{longtable,booktabs,array}
\usepackage{calc} % for calculating minipage widths
% Correct order of tables after \paragraph or \subparagraph
\usepackage{etoolbox}
\makeatletter
\patchcmd\longtable{\par}{\if@noskipsec\mbox{}\fi\par}{}{}
\makeatother
% Allow footnotes in longtable head/foot
\IfFileExists{footnotehyper.sty}{\usepackage{footnotehyper}}{\usepackage{footnote}}
\makesavenoteenv{longtable}
\usepackage{graphicx}
\makeatletter
\def\maxwidth{\ifdim\Gin@nat@width>\linewidth\linewidth\else\Gin@nat@width\fi}
\def\maxheight{\ifdim\Gin@nat@height>\textheight\textheight\else\Gin@nat@height\fi}
\makeatother
% Scale images if necessary, so that they will not overflow the page
% margins by default, and it is still possible to overwrite the defaults
% using explicit options in \includegraphics[width, height, ...]{}
\setkeys{Gin}{width=\maxwidth,height=\maxheight,keepaspectratio}
% Set default figure placement to htbp
\makeatletter
\def\fps@figure{htbp}
\makeatother

\KOMAoption{captions}{tableheading}
\makeatletter
\@ifpackageloaded{caption}{}{\usepackage{caption}}
\AtBeginDocument{%
\ifdefined\contentsname
  \renewcommand*\contentsname{Table of contents}
\else
  \newcommand\contentsname{Table of contents}
\fi
\ifdefined\listfigurename
  \renewcommand*\listfigurename{List of Figures}
\else
  \newcommand\listfigurename{List of Figures}
\fi
\ifdefined\listtablename
  \renewcommand*\listtablename{List of Tables}
\else
  \newcommand\listtablename{List of Tables}
\fi
\ifdefined\figurename
  \renewcommand*\figurename{Figure}
\else
  \newcommand\figurename{Figure}
\fi
\ifdefined\tablename
  \renewcommand*\tablename{Table}
\else
  \newcommand\tablename{Table}
\fi
}
\@ifpackageloaded{float}{}{\usepackage{float}}
\floatstyle{ruled}
\@ifundefined{c@chapter}{\newfloat{codelisting}{h}{lop}}{\newfloat{codelisting}{h}{lop}[chapter]}
\floatname{codelisting}{Listing}
\newcommand*\listoflistings{\listof{codelisting}{List of Listings}}
\makeatother
\makeatletter
\makeatother
\makeatletter
\@ifpackageloaded{caption}{}{\usepackage{caption}}
\@ifpackageloaded{subcaption}{}{\usepackage{subcaption}}
\makeatother

\ifLuaTeX
  \usepackage{selnolig}  % disable illegal ligatures
\fi
\usepackage{bookmark}

\IfFileExists{xurl.sty}{\usepackage{xurl}}{} % add URL line breaks if available
\urlstyle{same} % disable monospaced font for URLs
\hypersetup{
  pdftitle={Künstliche Intelligenz},
  colorlinks=true,
  linkcolor={blue},
  filecolor={Maroon},
  citecolor={Blue},
  urlcolor={Blue},
  pdfcreator={LaTeX via pandoc}}


\title{Künstliche Intelligenz}
\usepackage{etoolbox}
\makeatletter
\providecommand{\subtitle}[1]{% add subtitle to \maketitle
  \apptocmd{\@title}{\par {\large #1 \par}}{}{}
}
\makeatother
\subtitle{Richtiger Einsatz in Didaktik und Pädagogik \emph{Dr.~Wolfgang
Spahn} Oct.2024}
\author{}
\date{}

\begin{document}
\maketitle


\subsection{Agenda:}\label{agenda}

\begin{itemize}
\tightlist
\item
  KI verändert geistiges Arbeiten
\item
  KI Chatbots und Tools im Schulumfeld

  \begin{itemize}
  \tightlist
  \item
    Realitäten/Kontrolle des KI Einsatzes
  \item
    Förderung durch KI Einsatz
  \end{itemize}
\item
  konstruktivistisch-kognitive KI Lernsysteme

  \begin{itemize}
  \tightlist
  \item
    AIDu
  \item
    Vision
  \end{itemize}
\item
  Diskussion
\end{itemize}

\subsection{\_}\label{section}

KI verändert geistiges Arbeiten

\subsection{Erste Untersuchungen über KI/Mensch
Zusammenarbeit}\label{erste-untersuchungen-uxfcber-kimensch-zusammenarbeit}

758 consultants

\includegraphics[width=0.5\textwidth,height=\textheight]{images/havardBusinessSchool.webp}

4867 software developers

\includegraphics[width=0.5\textwidth,height=\textheight]{images/High skilled work AI.webp}

sehen 25-40\% Zunahme der Produktivität: skales with detached (-),
delegating (+) or immersive usage (++)

\subsection{körperliche Arbeit}\label{kuxf6rperliche-arbeit}

\textbf{1900}

\includegraphics[width=0.75\textwidth,height=\textheight]{images/1900_landwirtschaf_1.webp}

\textbf{2024}

\includegraphics[width=1\textwidth,height=\textheight]{images/claas mähdrescher.webp}

als Model der zu erwartenden Transformation?

\subsection{\_}\label{section-1}

KI Chatbots im Schulumfeld

\subsection{Chatbots in der Bildung}\label{chatbots-in-der-bildung}

\includegraphics[width=1\textwidth,height=\textheight]{images/2847_cartoon_back_to_school_1.webp}
\href{https://www.prpeak.com/opinion/cartoon-homework-assignments-7489368}{Source}

``Of the top 20 education apps in the U.S. App Store, five are AI agents
that help students with their school assignments, according to data from
Data.ai on May 21.''

\subsection{NZZ 2024}\label{nzz-2024}

Ein Zürcher Maturand sagt {[}Bestnote Deutsch!{]}: «Ich habe keines
meiner Bücher gelesen. Das hat alles die KI für mich gemacht»

\includegraphics[width=0.5\textwidth,height=\textheight]{images/NZZ Deutsch Matur.webp}
\includegraphics[width=0.5\textwidth,height=\textheight]{images/NZZ Faust.webp}

\subsection{KI Bots im Schulumfeld}\label{ki-bots-im-schulumfeld}

Vergleichsstudie an 1000 Lernenden:

Prüfungsvorbereitung mit und ohne KI Chatbots

\begin{itemize}
\tightlist
\item
  Unterschiedliche Hausaufgaben zur Mathe Prüfungsvorbereitung

  \begin{itemize}
  \tightlist
  \item
    Gruppe 1: Basis - traditionell
  \item
    Gruppe 2: +48\% - ChatGPT
  \item
    Gruppe 3: +127\% - einfacher sokratischer GPT-Tutor (nicht adaptiv)
  \end{itemize}
\item
  Prüfung über den Lernstoff

  \begin{itemize}
  \tightlist
  \item
    \textbf{ChatGPT Gruppen sind schlechter} als die traditionelle
    Gruppe
  \item
    GTP-Tutor Gruppe erzielt gleiche Ergebnisse wie die traditionelle
    Gruppe
  \end{itemize}
\end{itemize}

\includegraphics[width=0.7\textwidth,height=\textheight]{images/AI harm learning.webp}

\subsection{KI Bots im Schulumfeld}\label{ki-bots-im-schulumfeld-1}

Verbreitung und Akzeptanz generativer KI an Schulen und Hochschulen
(Bayern)

\begin{itemize}
\item
  Ergebnisse (Zusammenfassung
  \href{https://www.bidt.digital/publikation/verbreitung-und-akzeptanz-generativer-ki-an-schulen-und-hochschulen/}{\textgreater/})

  \begin{itemize}
  \tightlist
  \item
    \emph{\ldots{} haben unter den untersuchten Schülerinnen und
    Schülern 73 \%, unter den betrachteten Studierenden 78 \% bereits
    generative KI für die Text- und/oder Bildgenerierung genutzt.}
  \item
    \emph{\ldots, dass generative KI ihnen dabei geholfen hat, bessere
    Noten zu erhalten, ohne dafür eine angemessene Leistung erbracht zu
    haben, stimmen 42 \% der entsprechenden Schülerinnen und Schüler und
    45 \% der entsprechenden Studierenden eher oder voll und ganz zu.}
  \item
    \emph{Die Zufriedenheit mit der Nutzung textbasierter generativer KI
    in Schule oder Studium ist hoch. Nur eine Minderheit der
    untersuchten Nutzerinnen und Nutzer berichtet das Gegenteil.}
  \item
    \emph{\ldots, Herausforderung besteht in einer adäquaten
    Leistungsbewertung von Lernenden in Zeiten generativer KI, \ldots{}
    . Darum sollte dem Wunsch vieler Befragten Rechnung getragen werden
    und Prüfungsformate an Schulen und Hochschulen angepasst werden.}
  \end{itemize}
\end{itemize}

\includegraphics[width=0.7\textwidth,height=\textheight]{images/bidt_Studie_Generative_KI-an_Schulen_und_Hochschulen.webp}

\subsection{Was ist das Problem?}\label{was-ist-das-problem}

\textbf{KI im Schulumfeld ist per se nicht förderlich}

\begin{itemize}
\tightlist
\item
  KI Technologie {\textbf{reduziert Durchdringung}}

  \begin{itemize}
  \tightlist
  \item
    verleitet zu {\textbf{Kompetenzsimulation}}
  \item
    {\textbf{reduziert die Notwendigkeit von Wissen}} durch schnelle
    Verfügbarkeit von Information
  \end{itemize}
\item
  {\textbf{Miss- und Fehlinformation}} sind leicht verfüg-/erzeugbar
\item
  ChatBots verführen zu {\textbf{Plagiat und Betrug}}
\item
  Führt zur {\textbf{fehlerhaften Selbsteinschätzung}} der eigenen
  Kompetenz
\end{itemize}

Aber KI ist Realität. Daher ist \textbf{Wegschauen ist keine Lösung}!

\begin{figure}[H]

{\centering \includegraphics[width=1\textwidth,height=\textheight]{images/consumer_education1.webp}

}

\caption{DALL·E 3: Consumer Education}

\end{figure}%

\subsection{KI bewusste Pädagogik}\label{ki-bewusste-puxe4dagogik}

\textbf{Es braucht Methodik um die Vorteile von KI pädagogisch zu
nutzen. KI: }

\begin{itemize}
\tightlist
\item
  bringt Kapazität für {\textbf{individuelle Förderung}}
\item
  ermöglicht Schöpfungstiefe durch {\textbf{komplexere
  Aufgabenstellungen}}
\item
  fördert Durchdringung durch direkten {\textbf{Start an der
  Kompetenzgrenze}}
\item
  zielt auf tiefes Wissen durch {\textbf{Vernetzung von Wissen und
  Kompetenzen}}
\item
  motiviert durch lebensnahes, {\textbf{exploratives Lernen}} von
  komplexen Zusammenhängen
\end{itemize}

\begin{figure}[H]

{\centering \includegraphics[width=1\textwidth,height=\textheight]{images/explorer_education1.webp}

}

\caption{DALL·E 3: Education as Exploration}

\end{figure}%

\textbf{Das bedeutet eine enge Lernbegleitung durch Lehrpersonen und
Tools}

\subsection{KI bewusste Pädagogik}\label{ki-bewusste-puxe4dagogik-1}

\begin{itemize}
\tightlist
\item
  \textbf{Lernende}

  \begin{itemize}
  \tightlist
  \item
    entwickelt Motivation hin zu {\textbf{verankertem Wissen}}
  \item
    erlernt {\textbf{Sprach- und Dialogkompetenz}} um auf Augenhöhe mit
    KI zu kommunizieren
  \item
    arbeitet {\textbf{mit KI offen und selbstwirksam}}
  \item
    {\textbf{orchestriert KI}} Sitzungen um ein gewünschtes Projektziel
    zu erreichen
  \item
    erwirbt vor allem {\textbf{Basalwissen}}, das effektiven KI
    gestützten Einstieg in eine neue Domäne ermöglicht
  \end{itemize}
\item
  \textbf{Lehrende}

  \begin{itemize}
  \tightlist
  \item
    begeht mit den Lernenden {\textbf{individuelle Lernwege}}
  \item
    ermöglicht den Lernenden {\textbf{echte Kompetenzerfahrungen}}
  \item
    steigert {\textbf{Sprach/Dialogkompetenz}} der Lernenden
  \item
    lehrt {\textbf{Umgang mit Wissensüberfluss}}
  \item
    hat dafür {\textbf{Rüstzeug}}
  \end{itemize}
\end{itemize}

\begin{figure}[H]

{\centering \includegraphics[width=1\textwidth,height=\textheight]{images/Bergführerin.webp}

}

\caption{DALL·E 3: Seilschaft}

\end{figure}%

Forschungsfragen: Wie reduzieren auf {\textbf{Basalwissen}}? Welches
{\textbf{Rüstzeug}}?

\subsection{\_}\label{section-2}

KI bewusste Pädagogik

Intelligente Lernbegleitsysteme

\subsection{AIDu: ITS Forschung}\label{aidu-its-forschung}

\includegraphics[width=0.9\textwidth,height=\textheight]{images/AIdu.webp}
\includegraphics[width=0.4\textwidth,height=\textheight]{images/AIDu Welcome.webp}
\includegraphics[width=0.4\textwidth,height=\textheight]{images/AIDu Isaac.webp}

\subsection{Vision Bildung 2035}\label{vision-bildung-2035}

\textbf{Bildung unterstützt/ermöglicht durch KI durch}

\begin{itemize}
\tightlist
\item
  \textbf{Dialog} as Grundelement
\item
  \textbf{personalisierte Ansätze} für

  \begin{itemize}
  \tightlist
  \item
    individuelle Entwicklung
  \item
    eigenes Tempo
    \includegraphics[width=0.7\textwidth,height=\textheight]{images/student-study-group.webp}
  \end{itemize}
\item
  \textbf{Lernen in der Gruppe}

  \begin{itemize}
  \tightlist
  \item
    mit Lehrpersonen als Lernpartner
  \item
    mittels einem eigenverantwortlichen Lernprozess
  \end{itemize}
\end{itemize}

\includegraphics{index_files/mediabag/images/left-right-bold-arrow-icon.pdf}

\includegraphics[width=0.4\textwidth,height=\textheight]{images/Holzwerkstatt_2-1-scaled-crop.webp}
\includegraphics[width=0.4\textwidth,height=\textheight]{images/csm_Finn_mit_Sonde_ea9ce59730.webp}
\includegraphics[width=0.4\textwidth,height=\textheight]{images/schüler im labor.webp}

Sources:
\href{https://www.badische-zeitung.de/seit-30-jahren-wird-im-novartis-schullabor-in-basel-neugier-auf-wissenschaft-geweckt}{1}
\href{https://www.evim.de/news-veranstaltungen/news-und-veranstaltungen/detail-meldungen/hoehenflug-fuer-schueler-der-technik-werkstatt/}{2}
\href{https://freie-landschule.de/wp-content/uploads/2024/03/Holzwerkstatt_2-1-2048x1536.webp}{3}

\subsection{Zusammenfassung}\label{zusammenfassung}

\begin{itemize}
\tightlist
\item
  KI hat das Potential die Bildung zu revolutionieren

  \begin{itemize}
  \tightlist
  \item
    oder zu ruinieren
  \end{itemize}
\item
  KI Chatbots im Schulumfeld sind Realität

  \begin{itemize}
  \tightlist
  \item
    und werden heute (falsch?) genutzt
  \end{itemize}
\item
  KI bewusste Pädagogik

  \begin{itemize}
  \tightlist
  \item
    ist der Schlüssel
  \end{itemize}
\item
  KI Lernsysteme können dabei

  \begin{itemize}
  \tightlist
  \item
    die nötige Unterstützung bieten
  \end{itemize}
\end{itemize}

\begin{figure}[H]

{\centering \includegraphics[width=1\textwidth,height=\textheight]{images/explorer_education1.webp}

}

\caption{DALL·E 3: Education as Exploration}

\end{figure}%

\subsection{\_}\label{please_vote_id}

\emph{-} Danke für die Aufmerksamkeit \emph{-}
\href{https://wolfgangspahn.github.io/intro-IS2-KI-research.github.io}{\textgreater/}

\subsection{Dialoganalyse}\label{dialoganalyse}

\begin{itemize}
\tightlist
\item
  \textbf{Dialog A: High score}

  \begin{itemize}
  \tightlist
  \item
    AI: Let's solve this. Please read it
  \item
    Student: I have read it
  \item
    \ldots..
  \item
    Student: we get x
  \item
    AI: {So we bring everything with x on one side}
  \item
    Student: {8x = 10}
  \item
    AI: great you have done it, and now
  \item
    Student: I need to get x
  \item
    AI: {yes now you need to have just x on one side}
  \item
    Student: {x = 10/8}
  \item
    AI: nearly done
  \item
    Student: what do you mean
  \item
    AI: {you can simplify further}
  \item
    Student: {x = 5/4}
  \end{itemize}
\end{itemize}

\begin{itemize}
\tightlist
\item
  \textbf{Dialog B: Low score}

  \begin{itemize}
  \tightlist
  \item
    AI: Let's solve this. Please read it
  \item
    Student: I have read it
  \item
    AI: how would you solve this
  \item
    Student: {we need to isolate x to one side}
  \item
    AI: Good, just do it
  \item
    Student: {8x = 10}
  \item
    AI: And now
  \item
    Student: {x = 10/8 = 5/4}
  \end{itemize}
\item
  \textbf{Fragen}:

  \begin{itemize}
  \tightlist
  \item
    How much is the contribution of the student and AI in A and in B.
  \item
    What concepts are still not understood by the student in A
  \item
    What strategy towards solution does AI and student use in A and B
  \end{itemize}
\end{itemize}

\subsection{Qualitätskontrolle}\label{qualituxe4tskontrolle}

Da wir mit `open ended' Fragen arbeiten, ist die Qualitätskontrolle
schwierig.

\begin{itemize}
\tightlist
\item
  \textbf{Verwende die Vorteile von klassischer KI}

  \begin{itemize}
  \tightlist
  \item
    Konvertiere die Aufgabenstellung in ein computerverständliches
    Format
  \item
    nutzte symbolische KI Methoden zur Verifikation
  \end{itemize}
\item
  \textbf{Ansätze}

  \begin{itemize}
  \tightlist
  \item
    Computer Algebra Systeme für Mathematik, Physik, Chemie
  \item
    Knowledge Graphen für nicht mathematische Inhalte: Bio, Geschichte
    \ldots{} (?)
  \item
    Expertensysteme und Logik für Schlussfolgerungen (?)
  \end{itemize}
\end{itemize}

\includegraphics[width=1\textwidth,height=\textheight]{images/computer_algebra.webp}




\end{document}
